\documentclass[10pt,a4paper]{amsart}
\usepackage[margin=19mm]{geometry}
\usepackage{amsmath,amssymb,mathtools}
\usepackage[scr=rsfs,cal=euler]{mathalfa}
\usepackage{xfrac}
\usepackage[unicode,hidelinks,bookmarks=false]{hyperref}
\newcommand{\ie}{i.e.\ }
\newcommand{\cf}{cf.\ }
\newcommand{\resp}{respectively\ }
\newcommand{\Subsection}{\subsection}
\providecommand{\nobreakdash}{\nobreak\hbox{-}}
\setlength{\emergencystretch}{2em}
\multlinegap0pt
\usepackage{amssymb}
\usepackage{xcolor}
\usepackage{algorithm}\usepackage{algpseudocode}
\newtheorem{theorem}{Theorem}
\newcommand{\Sp}{\mathbb S}
\title{Improved bounds for the double cap conjecture}
\author{Domonkos Czifra, \'Akos D\'ucz, M\'at\'e Matolcsi, D\'aniel Varga, P\'al Zs\'amboki}
\date{Source: arXiv:2605.28709v1 (CC BY 4.0)}
\begin{document}
\maketitle

\noindent\emph{Source and licence.} Improved bounds for the double cap conjecture. Version arXiv:2605.28709v1 (CC BY 4.0), distributed by arXiv under the Creative Commons Attribution 4.0 International licence, http://creativecommons.org/licenses/by/4.0/, as recorded in the arXiv licence metadata for this exact version and shown on the abstract page https://arxiv.org/abs/2605.28709v1. Full licence notice: \texttt{licenses/arxiv-2605.28709-CC-BY-4.0.txt}.\par\smallskip

\noindent\textit{Editorial scope.} Counted unit: the article's theorem bound (source0.tex lines 297--307). The article's complete original proof of it is delivered with it (source0.tex lines 308--314, one \texttt{proof} environment of 7 lines). No mathematics is written, repaired or re-derived: the statement and the proof are the article's own lines, in the article's own notation, and no retained line is substituted or re-flowed.  Retained uncounted as the source-local context the proof needs: lines 157--159; lines 174--176; lines 234--236; lines 258--260; lines 262--264; lines 268--270.  Nothing was omitted from any retained span, so the package carries no omission record.  No retained line contains a citation, so the wrapper carries no bibliography and no bibliography entry.  No retained line contains a figure environment, an image-inclusion command or drawing code, so no image is delivered and the package declares no asset dependency.  Editorial formatting only, in addition: an AMS \texttt{amsart} preamble that restates the article's own declarations and macros used by the retained lines, namely the environments \texttt{theorem} on separate counters, together with the standard packages those lines need.  The retained environments print in this document as Theorem 1 in order of appearance, whereas the article prints the same items with the numbering of its own full document.  The complete native source of the article as distributed in the arXiv e-print for this exact version is bundled unchanged as \texttt{sources/201560/source0.tex}.
\begin{equation}\label{iet}
\sum_I a_X(I)= 1,    
\end{equation}

\begin{equation}\label{iec}
\sum_{I\supset Y} a_X(I)-\sum_{I\supset Z} a_X(I)=0  \ \ {\textrm{whenever \ }} Y\cong Z.
\end{equation}

\begin{equation}\label{pdf}
f_A(t)=\sum_{k=0}^\infty c_kP_k(t),  
\end{equation}

\begin{equation}\label{t1}
\sum_{2|k, k=0}^\infty c_k=\delta(A)  \ {\textrm{due to}} \ t=1 \ {\textrm{in}} \ \eqref{pdf}   
\end{equation}

\begin{equation}\label{t0}
\sum_{2|k, k=0}^\infty c_kP_k(0)=0  \ {\textrm{due to}} \ t=0 \ {\textrm{in}} \ \eqref{pdf},  
\end{equation}

\begin{equation}\label{axck}
\sum_{2|k, k=0}^\infty c_kP_k(\langle v_i, v_j \rangle) - \sum_{v_i,v_j\in I} a_X(I) =0 \ {\textrm{for all}} \ 1\le i\le j\le N.
\end{equation}

\begin{theorem}\label{bound}
Assume $X=\{x_1, -x_1, \dots, x_N, -x_N\}\subset \Sp^{n-1}$ is a finite, symmetric subset of the sphere, $v_i=\{x_i, -x_i\}$, $G_X$ is the orthogonality graph corresponding to $X$, and real parameters $q_1$, $s_{Y,Z}$ (for subsets $Y$ and $Z$ of $G_X$, where $Y\cong Z$), $w_0$ and $w_{i,j}$ (for all $1\le i\le j\le N$)  are given such that the inequalities 
\begin{equation}\label{forck}
w_0P_k(0) + \sum_{i\le j}w_{i,j}P_k(\langle v_i, v_j \rangle )\ge 1 \ {\textrm{for all}} \ k\ge 0, \ k \  {\textrm{is even}} 
\end{equation}

\begin{equation}\label{forax}
q_1 - \sum_{v_i, v_j\in I}w_{i,j} + \sum_{I\subset Y}s_{Y,Z} -\sum_{I\subset Z}s_{Y,Z}\ge 0 \ {\textrm{for all}} \ I\subset X, I \  {\textrm{is independent}}
\end{equation}
are satisfied. Then $\alpha_n\le q_1$.    
\end{theorem}

\begin{proof}
The proof is a direct application of linear duality, as follows. Let us use the notation 
$\beta_k=w_0P_k(0)+\sum_{i\le j}w_{i,j}P_k(\langle v_i, v_j \rangle )$, and $\gamma_I=q_1 - \sum_{v_i, v_j\in I}w_{i,j} + \sum_{I\subset Y}s_{Y,Z} -\sum_{I\subset Z}s_{Y,Z}$. Consider the weighted sum of equations \eqref{iet}, \eqref{iec}, \eqref{t0} and \eqref{axck} with respective weights $q_1$, $s_{Y,Z}$, $w_0$ and $w_{i,j}$. This weighted sum gives us $\sum_{2|k, k=0}^\infty \beta_k c_k +\sum_{I\subset X}\gamma_Ia_X(I)= q_1$. Therefore, using the assumptions $\beta_k\ge 1, \gamma_I\ge 0$, equation \eqref{t1}, and the nonnegativity of the variables $c_k$ and  $a_X(I)$ we obtain 
\begin{equation}\label{last}
\delta(A)=\sum_{2|k, k=0}^\infty  c_k \le \sum_{2|k, k=0}^\infty \beta_k c_k +\sum_{I\subset X}\gamma_Ia_X(I) =q_1.       
\end{equation}
\end{proof}

\end{document}
